% page 606 of the facsimile (image p-605 of the scan)
% block 147.3 x 263.6 mm ; left margin 8.0 mm
\pgc{-12.700mm}{0.000mm}{
\l{}
\l{}
\l{}
\l{}
\l{\cel{3}598.}
\l{}
\l{\sou{trulo}. (\sou{tekn.})\cel{2}Utensilo spatulo-forma por dilutar ed}
\l{\cel{3}aplikar gipso, mortero, cemento, e c.- EFS.}
\l{}
\l{\sou{trumpo}. (\sou{karto-ludo}). La emblemo qua preponderas omna ce-}
\l{\cel{3}terai, e quan onu determinas singlafoye, sive per}
\l{\cel{3}vidigar (lor la distributo kun nur la reverso videbla)}
\l{\cel{3}la averso di la distributo-lasta - sive per konvenciono}
\l{\cel{3}irgaltra. - DE.}
\l{}
\l{\sou{trumpeto}. (\sou{muziko}). Sufl-instrumento ek latuno od altra}
\l{\cel{3}metalo, di qua la sono esas tridanta, boranta - uzata}
\l{\cel{3}por la ario-soni armeala, por la proklami sur placi.}
\l{\cel{3}- DEFIS.}
\l{}
\l{\sou{trunko}. I. Korpo di arboro, sen la branchi nek la radiki.}
\l{\cel{3}- II. (\sou{geom}.)Segmento di korpo definita. (Ex.: prismo-}
\l{\cel{3}trunko, kono-trunko, e c.) - EFIS.}
\l{}
\l{\sou{trunkono}. I. (\sou{tekn.}) Elemento di la vapor-kaldieri, e di}
\l{\cel{3}la tanki e cisterni metala, cilindra, qua konsistas ye}
\l{\cel{3}ringego de fer-tolo. - II. (\sou{arkitekt.})\cel{2}Parto infra di}
\l{\cel{3}fusto. -}
\l{}
\l{\sou{trupo}. Asemblitaro ek ula quanto de personi qui kuniras,}
\l{\cel{3}agas ensemble, interkonsente. - II. Asemblitaro ek}
\l{\cel{3}animali domestika quan onu edukas, nutras, duktas a}
\l{\cel{3}la pastureyo kom ensemblo. - DEF.}
\l{}
\l{\sou{trusar}. (\sou{trans.}) Eskartar de sulo per faldi (kozo qua}
\l{\cel{3}pendas). - F.}
\l{}
\l{\sou{trusto}. (\sou{komerco}). Sindikato de spekulisti, fondita por}
\l{\cel{3}plugrandigar la preco di varo o di acioni, per akapa-}
\l{\cel{3}rar li por pose vendar li granda-preco. - DEFIS.}
\l{}
\l{\sou{truto}. (\sou{zool.})\cel{2}Fisho ek la genero \textquotedbl{}salmoni\textquotedbl{}, kun pelo}
\l{\cel{3}makuletoza, di qua la karno esas tre delikate sapo-}
\l{\cel{3}roza. - L.\cel{1}\sou{salmo fario}. - EFIS.}
\l{}
\l{\sou{truvero}. (\sou{historio literat.}) Poeto qua kompozis per la}
\l{\cel{3}linguo olima di la nordo-regioni di Francia. - DEFIS.}
\l{}
\l{\sou{tu.}\cel{3}(\sou{gram.})\cel{2}Pronomo personala, vokativo, uzata kande}
\l{\cel{3}on parolas o skribas a persono kun qua onu relatas pa-}
\l{\cel{3}sable intime (pro la ne-ceremonialeso di ta pronomo);}
\l{\cel{3}a persono kun qua onu relatas intime e mem tre intime ;}
\l{\cel{3}a persono qua esas tre yuna - ed a persono quan onu}
\l{\cel{3}desprizas. - DEFIRS.}
\l{}
\l{\sou{tualetar.}\cel{1}(\sou{netrans.}) Lavar, razar, kuafar, vestizar ed}
\l{\cel{3}ornar su, generale en kabineto. - DEFIRS.}
\l{}
\l{\sou{tubo}. (\sou{tekn.}) Cilindro kava, naturala od artificala,quan}
\l{\cel{3}onu uzas por pasigar, ula-direcione, liquidi, fluidi. -}
\l{\cel{3}EFIS.}
}
